\section{Information, comparisons and retained results}
\label{sec:comparison}
\subsection{Boundary distance, lens data and obstacle travelling times}
The exact equivalence in Proposition~\ref{prop:lens} identifies the
information category of the local result. On a physical active rectangle,
two endpoint density functions determine a travel-time generating function
and a volume normalization; conversely these two objects determine the
densities. The affine elimination is elementary. The geometric content
lies in recovering two unknown reflecting arcs from the diagonal Hessian,
and the new global content lies in recovering incidence and the unmarked
period lattice from an unlabelled collection. We do not claim that endpoint
travel times, their tangential derivatives, or the passage to a scattering
relation are new inverse-geometric concepts.

Stefanov, Uhlmann and Vasy~\cite{SUV}, Theorem~1.1, recover a Riemannian
metric locally near a strictly convex boundary point from boundary-distance
data in dimension at least three, modulo a boundary-fixing diffeomorphism.
Their Theorem~1.3 gives global lens rigidity under a convex-foliation
hypothesis. The unknown there is a metric on a manifold with identified
boundary. Here the ambient metric is Euclidean and known, while both
reflecting curves are unknown; the measured coordinates are intrinsic
arclengths on a selected unknown source obstacle. The diagonal action
is $W(s,s)=2\ell(s)>0$, rather than the zero diagonal of an ordinary
boundary-distance function. The Schur complement in Lemma~\ref{lem:curvature}
uses a stationary interior reflection. Thus the stated metric-rigidity
theorems do not directly give our two-arc formulas; nor does our selected
planar return experiment imply their much more general metric recovery.

Gurfinkel, Noakes and Stoyanov~\cite{GNS} consider travelling times between
points on a known enclosing boundary in a curved ambient space. Their
Theorem~1 establishes conjugacy of the nontrapping scattering flows from
almost-equal travelling times. Their Theorem~2 determines finite unions of
strictly convex obstacles in a two-dimensional nonpositively curved
manifold from those travelling times. Their Lemma~7 makes explicit how
endpoint derivatives recover the incident and outgoing directions.
Noakes and Stoyanov~\cite{NS}, Theorem~1 and Sections~3--4, likewise relate
travelling-time data to conjugacy and give a constructive reconstruction
for two convex planar obstacles. These are close comparisons, not merely
background spectral references. They use families of externally observed
scattering trajectories; our data select a local source--opposite--source
return and record intrinsic endpoints on an unknown reflecting component.
The local origins and selected branch remain marks. No theorem here
converts our finite channel collection into their full exterior data.

The volume normalization has a similarly established context. Stoyanov's
extension of Santal\'o's formula~\cite{Santalo} relates phase volume and
travelling-time averages and applies to obstacle-volume recovery in the
presence of trapping. Our identity for $A$ is instead a local consequence
of the unconditional residual-time density. Its use in
$V=A+\sum_i\area(C_i)$ converts geometric cycle displacements into
an integer period index. It is this combination, not a general claim
of first recovering volume from dynamics, that enters the descent theorem.

Marked-length results have a different information set. De Simoi, Kaloshin
and Leguil~\cite{DSKL} determine analytic dispersing configurations under
axial symmetry and genericity assumptions from marked lengths. Finamore
and Leguil~\cite{FL} prove isometric rigidity of finite-horizon Sinai
billiards from an enriched marked length spectrum. We neither identify
our endpoint laws with either spectrum nor infer a stronger spectral
rigidity theorem. The present no-finite-horizon assertion concerns the
selected local channels. The generic shape-separation and primitive-cycle
hypotheses are conditions of our unlabelled reconstruction, not conclusions
about arbitrary dispersing tables. These comparisons delimit the results;
they are not an exhaustive priority classification.

\subsection{What is removed, and what is retained}
Theorem~\ref{thm:descent-main} removes the supplied obstacle identifications,
incidence graph, cross-channel contact offsets, relative edge signs and
integer lattice-copy labels on the stated shape-separated class. The
record itself still pairs its two times, specifies one local return
orientation, and uses a marked contact origin and a coherent arclength
reversal. The finite theorem also uses onset brackets and uniform priors.
``Unlabelled'' therefore describes the relations \emph{between} records,
not the absence of all experimental selection.

The area calibration is indispensable to the arithmetic reduction.
Without it, containing two independent geometric periods does not bound
the index of a finer lattice. With it, the finite list is exact, and a
primitive determinant is a directly verifiable sufficient certificate
of uniqueness. Proposition~\ref{prop:ambiguity} shows that even the raw
normalization and rank-two geometry need not give uniqueness at larger
index. This obstruction persists for exact physical density functions;
it is not a statistical artifact or an inference from formal series.

The local Lipschitz estimate is conditional on the physical bounds in
Proposition~\ref{prop:local-stable}. The global exponent is additionally
conditional on the analytic strip and continuation geometry, as in
\cite{Trefethen}, and the separation margins in Section~\ref{sec:stability}.
The integer-locking argument is a comparison of physical candidates with
known integrality, not the formation of an integer span from noisy real
vectors. None of these estimates establishes a minimax exponent.

\subsection{Comparison with the two retained regularizations}
The complete preceding article is Supplement R. Its moment-compressed
experiment uses masses, signed first moments and quadratic moments.
For a contact-jet order $K$ its whole-table bound is
\[
 C\left(A_K\epsilon^{1/(\lceil K/2\rceil+2)}+Bq^K\right)^{\vartheta},
                       \qquad 0<q,\vartheta<1.
\]
The constants $A_K$ are finite at each order but need not be uniformly
bounded. A slowly increasing regularization order gives consistency.
That theorem remains useful when the acquisition stores moment means
rather than full endpoint histograms; the exact-germ version also exposes
the even and odd triangular blocks coefficient by coefficient.

Supplement R's full-law experiment instead gives
$C(h^\beta+\epsilon h^{-4})^\vartheta$ for absolute cell-probability
accuracy $\epsilon$. It has no growing jet order, but each of its two
windows is a growing two-dimensional histogram. The present paper retains
that deterministic local estimate. Using the bound $p_B\le Ch^2$ on
cell probabilities improves the preparation accounting to
\eqref{eq:rate-main}. Simultaneously, shape separation and period-index
locking permit recovery of global combinatorial data previously supplied.
The power in the new phase-preparation bound and the power in the old
absolute-mean-error bound refer to different error parameters. They are
not competing minimax exponents on a common unqualified sensor model.

The full-law route is the primary article because it reconstructs arcs
as functions before continuation and supports the unlabelled descent.
The complete moment inverse, its dimension-sharp finite-jet designs,
registered rigidity for symmetric tables, and the old order-dependent
regularization are supplied unchanged in Supplement R. Supplement S
retains the smooth relative boundary law and its Volterra and Abel
proofs. This separation changes the reading order, not the mathematical
content or status of those results.

Finally, the retained count-data fibers are finite, formal, or smooth
Taylor-series statements in their specified categories. Neither the
full-law inverse nor the period ambiguity theorem implies an exact
analytic count-only rigidity or nonrigidity theorem. In particular,
equality of all smooth Taylor coefficients is not used to identify
smooth functions, and a formal inverse series is not assumed to converge.
